% Options for packages loaded elsewhere
\PassOptionsToPackage{unicode}{hyperref}
\PassOptionsToPackage{hyphens}{url}
\PassOptionsToPackage{dvipsnames,svgnames,x11names}{xcolor}
\documentclass[
  11pt,
]{article}
\usepackage{xcolor}
\usepackage[margin=1in]{geometry}
\usepackage{amsmath,amssymb}
\setcounter{secnumdepth}{-\maxdimen} % remove section numbering
\usepackage{iftex}
\ifPDFTeX
  \usepackage[T1]{fontenc}
  \usepackage[utf8]{inputenc}
  \usepackage{textcomp} % provide euro and other symbols
\else % if luatex or xetex
  \usepackage{unicode-math} % this also loads fontspec
  \defaultfontfeatures{Scale=MatchLowercase}
  \defaultfontfeatures[\rmfamily]{Ligatures=TeX,Scale=1}
\fi
\usepackage{lmodern}
\ifPDFTeX\else
  % xetex/luatex font selection
  \setmonofont[]{DejaVu Sans Mono}
\fi
% Use upquote if available, for straight quotes in verbatim environments
\IfFileExists{upquote.sty}{\usepackage{upquote}}{}
\IfFileExists{microtype.sty}{% use microtype if available
  \usepackage[]{microtype}
  \UseMicrotypeSet[protrusion]{basicmath} % disable protrusion for tt fonts
}{}
\makeatletter
\@ifundefined{KOMAClassName}{% if non-KOMA class
  \IfFileExists{parskip.sty}{%
    \usepackage{parskip}
  }{% else
    \setlength{\parindent}{0pt}
    \setlength{\parskip}{6pt plus 2pt minus 1pt}}
}{% if KOMA class
  \KOMAoptions{parskip=half}}
\makeatother
\usepackage{color}
\usepackage{fancyvrb}
\newcommand{\VerbBar}{|}
\newcommand{\VERB}{\Verb[commandchars=\\\{\}]}
\DefineVerbatimEnvironment{Highlighting}{Verbatim}{commandchars=\\\{\}}
% Add ',fontsize=\small' for more characters per line
\newenvironment{Shaded}{}{}
\newcommand{\AlertTok}[1]{\textcolor[rgb]{1.00,0.00,0.00}{\textbf{#1}}}
\newcommand{\AnnotationTok}[1]{\textcolor[rgb]{0.38,0.63,0.69}{\textbf{\textit{#1}}}}
\newcommand{\AttributeTok}[1]{\textcolor[rgb]{0.49,0.56,0.16}{#1}}
\newcommand{\BaseNTok}[1]{\textcolor[rgb]{0.25,0.63,0.44}{#1}}
\newcommand{\BuiltInTok}[1]{\textcolor[rgb]{0.00,0.50,0.00}{#1}}
\newcommand{\CharTok}[1]{\textcolor[rgb]{0.25,0.44,0.63}{#1}}
\newcommand{\CommentTok}[1]{\textcolor[rgb]{0.38,0.63,0.69}{\textit{#1}}}
\newcommand{\CommentVarTok}[1]{\textcolor[rgb]{0.38,0.63,0.69}{\textbf{\textit{#1}}}}
\newcommand{\ConstantTok}[1]{\textcolor[rgb]{0.53,0.00,0.00}{#1}}
\newcommand{\ControlFlowTok}[1]{\textcolor[rgb]{0.00,0.44,0.13}{\textbf{#1}}}
\newcommand{\DataTypeTok}[1]{\textcolor[rgb]{0.56,0.13,0.00}{#1}}
\newcommand{\DecValTok}[1]{\textcolor[rgb]{0.25,0.63,0.44}{#1}}
\newcommand{\DocumentationTok}[1]{\textcolor[rgb]{0.73,0.13,0.13}{\textit{#1}}}
\newcommand{\ErrorTok}[1]{\textcolor[rgb]{1.00,0.00,0.00}{\textbf{#1}}}
\newcommand{\ExtensionTok}[1]{#1}
\newcommand{\FloatTok}[1]{\textcolor[rgb]{0.25,0.63,0.44}{#1}}
\newcommand{\FunctionTok}[1]{\textcolor[rgb]{0.02,0.16,0.49}{#1}}
\newcommand{\ImportTok}[1]{\textcolor[rgb]{0.00,0.50,0.00}{\textbf{#1}}}
\newcommand{\InformationTok}[1]{\textcolor[rgb]{0.38,0.63,0.69}{\textbf{\textit{#1}}}}
\newcommand{\KeywordTok}[1]{\textcolor[rgb]{0.00,0.44,0.13}{\textbf{#1}}}
\newcommand{\NormalTok}[1]{#1}
\newcommand{\OperatorTok}[1]{\textcolor[rgb]{0.40,0.40,0.40}{#1}}
\newcommand{\OtherTok}[1]{\textcolor[rgb]{0.00,0.44,0.13}{#1}}
\newcommand{\PreprocessorTok}[1]{\textcolor[rgb]{0.74,0.48,0.00}{#1}}
\newcommand{\RegionMarkerTok}[1]{#1}
\newcommand{\SpecialCharTok}[1]{\textcolor[rgb]{0.25,0.44,0.63}{#1}}
\newcommand{\SpecialStringTok}[1]{\textcolor[rgb]{0.73,0.40,0.53}{#1}}
\newcommand{\StringTok}[1]{\textcolor[rgb]{0.25,0.44,0.63}{#1}}
\newcommand{\VariableTok}[1]{\textcolor[rgb]{0.10,0.09,0.49}{#1}}
\newcommand{\VerbatimStringTok}[1]{\textcolor[rgb]{0.25,0.44,0.63}{#1}}
\newcommand{\WarningTok}[1]{\textcolor[rgb]{0.38,0.63,0.69}{\textbf{\textit{#1}}}}
\usepackage{longtable,booktabs,array}
\usepackage{calc} % for calculating minipage widths
% Correct order of tables after \paragraph or \subparagraph
\usepackage{etoolbox}
\makeatletter
\patchcmd\longtable{\par}{\if@noskipsec\mbox{}\fi\par}{}{}
\makeatother
% Allow footnotes in longtable head/foot
\IfFileExists{footnotehyper.sty}{\usepackage{footnotehyper}}{\usepackage{footnote}}
\makesavenoteenv{longtable}
\setlength{\emergencystretch}{3em} % prevent overfull lines
\providecommand{\tightlist}{%
  \setlength{\itemsep}{0pt}\setlength{\parskip}{0pt}}
\usepackage{bookmark}
\IfFileExists{xurl.sty}{\usepackage{xurl}}{} % add URL line breaks if available
\urlstyle{same}
\hypersetup{
  pdftitle={Warden: Event-Gated Threshold Conditional-Decryption for Guaranteed-but-Deferred Encrypted Delivery},
  pdfauthor={Sandeep Nandal},
  colorlinks=true,
  linkcolor={blue},
  filecolor={Maroon},
  citecolor={Blue},
  urlcolor={blue},
  pdfcreator={LaTeX via pandoc}}

\title{Warden: Event-Gated Threshold Conditional-Decryption for
Guaranteed-but-Deferred Encrypted Delivery}
\author{Sandeep
Nandal\thanks{BytesBrains Pte Ltd, \texttt{sandeep@nandal.in}. Solo authorship reflects the work as done; a co-author position is open to an academic collaborator who contributes to the formal peer-reviewed version.}}
\date{Draft v1 · June 2026}

\begin{document}
\maketitle

\textbf{Status disclosure (read first).} Warden is \textbf{deployed on a
public testnet (Base Sepolia) and is pre-audit.} This paper presents a
\emph{design, threat model, and early deployment experience}, not an
audited production system. The current testnet federation --- five
nodes, but all operated by a single party (BytesBrains, on one
infrastructure provider) --- therefore has, by construction, \textbf{no
timing-security} (see §6, W4); it exercises the protocol, integration,
and liveness, not the trust model. Independent operators are actively
being recruited (§9). We state this plainly because, for the audience
this work addresses, calibrated honesty about limitations is the
contribution's credibility, not a footnote to it.

\begin{center}\rule{0.5\linewidth}{0.5pt}\end{center}

\section{Abstract}\label{abstract}

Time-lock encryption lets a sender make a ciphertext readable only after
a predetermined time, and recent constructions such as \emph{tlock}
realize it efficiently over a threshold-BLS randomness beacon (drand).
But an important class of applications needs a \emph{conditional} rather
than a \emph{temporal} trigger, over a \emph{decades-long} horizon, with
the data author \emph{absent} at release time: a dead-man's-switch that
delivers a sealed message to designated recipients only if its owner
stops checking in --- for digital estates, personal-safety triggers, and
``your work survives you'' press-freedom scenarios.

We present \textbf{Warden}, an event-gated threshold
conditional-decryption network. A federation of independent operators
holds shares of an identity-based master key (Boneh--Franklin IBE over
threshold BLS, established by distributed key generation) and releases a
per-message partial decryption key only when an on-chain condition --- a
resettable smart-contract flag --- is observed true at finalized block
depth. Payloads are \emph{double-wrapped}: an outer condition-gate the
federation can open, and an inner recipient-key layer it cannot, so
operators learn release \emph{timing} but never \emph{content}.

We formalize the resettable-condition trigger and the owner-absent,
decades-horizon threat model; describe master-key permanence across
operator churn via resharing; give an honest security analysis
(threshold timing guarantees, finality/reorg leakage, collusion and
traitor-tracing); and report a Rust implementation deployed as a
five-node testnet federation that runs end-to-end on Base Sepolia,
integrated with the live Maktub protocol and a mobile client --- a
single-trust-domain deployment that exercises the protocol and its
integration, with independent-operator security as the explicit next
step. Warden reuses an audited cryptographic substrate (drand's
threshold-BLS/IBE/DKG stack) and contributes the \emph{system, threat
model, and conditional-trigger design} --- not new cryptography. It is
intended to run as token-free public-good infrastructure in the drand /
League-of-Entropy tradition. The precise contribution claims and
explicit non-claims are enumerated in §1.3 and §1.4.

\begin{center}\rule{0.5\linewidth}{0.5pt}\end{center}

\section{1. Introduction}\label{introduction}

\subsection{1.1 The gap}\label{the-gap}

Time-lock encryption (TLE) lets a sender seal a message that becomes
readable only after a chosen point in the future, without trusting any
single party to keep time. The deployed state of the art is
\textbf{tlock} {[}Gailly--Melissaris--Romailler 2023{]}, which observes
that drand's threshold-BLS randomness beacon is \emph{already} an
identity-based-encryption key-generator: the beacon's signature on a
future round number \emph{is} the IBE decryption key for that round. A
sender encrypts ``to round R''; once the League-of-Entropy federation
publishes round R's signature, anyone can decrypt --- offline, with no
interaction at decryption time.

This is powerful but fixes the trigger to a \textbf{clock}. A large and
practically important class of applications needs the release to be
gated not on \emph{time} but on a \emph{condition}, and over horizons
and absence assumptions that the timelock literature does not target:

\begin{itemize}
\tightlist
\item
  \textbf{Conditional, not temporal.} ``Release if and only if event E
  has happened,'' where E is some externally observable fact --- not
  ``release at time T.''
\item
  \textbf{Resettable.} The condition is not monotone-by-fiat from the
  sender's side: the owner actively keeps it \emph{false} by checking
  in, and it flips \emph{true} only when they stop. This is the
  dead-man's-switch shape, and it is the inverse of a timelock (where
  the countdown is unstoppable once started).
\item
  \textbf{Decades-horizon and owner-absent.} The ciphertext may sit
  dormant for years or decades, and --- definitionally --- the author is
  \emph{not present} at release time to participate, re-key, or
  intervene. Most timelock and threshold work assumes short horizons and
  a live participant.
\end{itemize}

\subsection{1.2 Motivating application: Maktub /
Veil}\label{motivating-application-maktub-veil}

Warden is motivated by \textbf{Maktub}, a decentralized
conditional-execution protocol on the Base L2 whose single primitive is:
\emph{if the owner does not check in within a chosen interval, deliver
an encrypted payload to designated recipients.} Maktub's on-chain
contract (\texttt{MaktubCore}, ``Beat'') is immutable and holds no
plaintext; it merely records, per \emph{beat}, whether the owner has
gone silent (an \texttt{executed} status the contract latches to
\texttt{true} once a timer lapses and an executor calls it).

The capability this enables for end users --- \emph{sealed until you go
silent; revocable while you are present; end-to-end encrypted to the
recipient} --- is called \textbf{Veil}. Veil's three named use cases are
equal: digital-estate inheritance (seed phrases, passwords, documents to
heirs), personal-safety triggers (hikers, field workers, solo
travelers), and press-freedom \textbf{``your work survives you''}
(material a journalist prepared is delivered to recipients they
designated if they become unable to deliver it themselves --- framed as
\emph{work survival}, e.g.~accident or illness, \textbf{never} as
adversarial-state anonymity).

Maktub's contract enforces the \emph{timer and the trigger}. It does
\textbf{not} by itself enforce the \emph{time-confidentiality} half of
Veil: nothing in a public smart contract can hold a payload
unreadable-until-trigger, because anything the contract can read, the
world can read. That half is exactly what Warden supplies.
\texttt{MaktubCore} requires \textbf{no modification} to support Warden
--- Warden only \emph{reads} the beat's status.

\subsection{1.3 Contribution}\label{contribution}

We claim the following, precisely:

\begin{enumerate}
\def\labelenumi{\arabic{enumi}.}
\tightlist
\item
  \textbf{A conditional, resettable trigger.} Release is gated on
  \emph{mutable on-chain contract state} (a flag the owner resets by
  checking in), evaluated at finalized block depth --- not a fixed time
  or round. This problem shape is under-explored relative to tlock and
  time-lock puzzles, and the \emph{combination} it sits in is not served
  by deployed conditional-decryption systems (§2.6).
\item
  \textbf{An owner-absent, decades-horizon threat model}, with
  master-key permanence maintained across operator churn by
  \emph{resharing} (the master public key --- and hence the
  decryptability of all dormant ciphertexts --- survives a completely
  changed operator set).
\item
  \textbf{A double-wrap construction} that cleanly separates
  \emph{timing control} (the threshold federation) from \emph{content
  confidentiality} (a recipient-held key, pure mathematics): operators
  can open the outer condition-gate but \textbf{never} see plaintext.
\item
  \textbf{An honest, extensive limitation analysis} (W1--W11):
  early-release via collusion, liveness/delay, finality/reorg leakage,
  operator/infra capture, master-key permanence as a double-edged sword,
  the non-post-quantum horizon, condition-evaluation trust tiers, the
  open-request metadata channel, cross-sender identity collision,
  share-key registration, and request-channel denial-of-service.
\item
  \textbf{An open-source (MIT) implementation and live testnet
  deployment:} a single Rust core driving a node daemon, a WASM web
  client, and an FFI mobile client, running end-to-end as a five-node
  federation on Base Sepolia and integrated with the Maktub protocol (a
  single-trust-domain testnet; independent operators are being recruited
  --- §9).
\end{enumerate}

The contribution is a \textbf{systems / applied} one. Gen-1 Warden
reuses drand's audited threshold-BLS + Boneh--Franklin IBE + DKG +
resharing essentially verbatim and replaces only the \emph{trigger}
(clock → on-chain condition) and the engineering around it. We are
explicit about this because aiming such work at cryptographic-theory
venues would be both wrong and counterproductive: there is no new
primitive here, and the genuinely novel surface (a resettable-condition
release gate with a finality policy) is a systems contribution.

\subsection{1.4 Explicit non-claims}\label{explicit-non-claims}

The same honesty discipline that governs all Maktub copy governs this
paper. We do \textbf{not} claim:

\begin{itemize}
\tightlist
\item
  That timing/revocation is \textbf{trustless}, \textbf{absolute}, or
  \textbf{mathematically guaranteed}. It rests on a \textbf{threshold
  honesty assumption} (no early release unless \texttt{t} operators
  collude; no denial unless fewer than \texttt{t} stay honest and live
  --- §6.1) \emph{and} on operators correctly observing the on-chain
  condition. The trustless ideal is witness encryption (§2.4), which is
  not deployable today; Warden is the best deployable approximation.
\item
  That Warden provides \textbf{metadata privacy or anonymity.} The
  existence of a sealed message, its release timing, its public
  condition, and (for Veil) its \texttt{beatId}, recipient set, and
  sender-recipient relationship are \textbf{public on-chain by design.}
  Warden hides \emph{content until the trigger} --- never \emph{that
  something exists}.
\item
  That gen-1 is \textbf{post-quantum.} It is pairing-based (BLS12-381);
  the realistic horizon for the threshold layer is on the order of
  years, and no deployable post-quantum threshold-IBE replacement
  exists. The envelope is algorithm-versioned so the cryptography can
  migrate (§10).
\end{itemize}

\subsection{1.5 Roadmap of the paper}\label{roadmap-of-the-paper}

§2 surveys the cryptographic background and positions Warden against
timelock, time-lock puzzles, witness encryption, and signature-based
witness encryption. §3 formalizes the resettable-condition trigger and
the threat model. §4 gives the system design (double-wrap envelope,
condition model, network planes, finality policy). §5 gives the
cryptographic construction (DKG, encryption, threshold release,
resharing). §6 is the honest security analysis (W1--W11) --- the
credibility core. §7 describes the implementation; §8 the testnet
deployment experience. §9 covers federation governance and
sustainability; §10 limitations and future work; §11 ethics and
responsible disclosure; §12 concludes.

\begin{center}\rule{0.5\linewidth}{0.5pt}\end{center}

\section{2. Background and Related
Work}\label{background-and-related-work}

\subsection{2.1 Threshold cryptography and distributed key
generation}\label{threshold-cryptography-and-distributed-key-generation}

A \emph{(t, n)}-threshold scheme splits a secret among \texttt{n}
parties such that any \texttt{t} can jointly use it but any \texttt{t−1}
learn nothing. Shamir secret sharing {[}Shamir 1979{]} is the
foundation; Feldman {[}1987{]} and Pedersen {[}1991{]} add
\emph{verifiability} (VSS), letting parties detect a cheating dealer.
\textbf{Distributed key generation} (DKG) {[}Gennaro et al.~1999{]}
removes the dealer entirely: \texttt{n} parties jointly produce a public
key and individual private-key shares such that the full private key is
\emph{never assembled by anyone}. drand uses Pedersen's DKG built from
parallel Feldman VSS, and contributes a hardened orchestration layer
around it (a durable state machine, explicit joiner/remainer/leaver
lifecycle).

\textbf{Threshold BLS} {[}Boneh--Lynn--Shacham 2001; Boldyreva 2003{]}
is the release primitive Warden inherits. A BLS signature on message
\texttt{m} under secret key \texttt{sk} is \texttt{sk·H(m)} for a
hash-to-curve \texttt{H}. Splitting \texttt{sk} into Shamir shares
\texttt{sk\_i}, each party emits a partial
\texttt{sig\_i\ =\ sk\_i·H(m)}; any \texttt{t} partials Lagrange-combine
\emph{in the exponent} into the full signature \texttt{sk·H(m)}.
Signatures are deterministic (hence unbiasable), which is what makes the
drand beacon usable as a key generator.

\subsection{2.2 Identity-based encryption and timelock
(tlock)}\label{identity-based-encryption-and-timelock-tlock}

Boneh--Franklin IBE {[}BF 2001{]} lets a sender encrypt to an arbitrary
\emph{identity string} using only a master public key; the matching
decryption key \texttt{msk·H1(id)} is generated later by the key
authority. The Boneh--Franklin ciphertext \texttt{(U,\ V,\ W)} with the
Fujisaki--Okamoto transform {[}FO 1999{]} achieves CCA security in the
random-oracle model.

\textbf{tlock} {[}Gailly--Melissaris--Romailler 2023{]} makes the
crucial observation: a threshold-BLS beacon \emph{is} an IBE private-key
generator that does not know it is one. Set the IBE identity to a future
round number; the beacon's BLS signature over that round
(\texttt{msk·H1(round)}) \emph{is exactly} the IBE decryption key for
that identity. No extra protocol is needed --- when the beacon naturally
publishes the round signature, the timelock opens. tlock requires the
beacon's \textbf{unchained} mode (the message signed for a future round
is predictable: \texttt{H(round)} alone, not chained on the previous
signature), and it packages the result as a hybrid age-encryption
envelope with a
\texttt{-\textgreater{}\ tlock\ \{round\}\ \{chainHash\}} stanza.

\textbf{The single most important fact Warden inherits from this line:}
nothing in the IBE math requires the identity to be a round number.
\emph{The identity is whatever string the network agrees to sign.}
Warden swaps \texttt{H(round)} for \texttt{H(condition)} and the
time-elapsed trigger for an on-chain-boolean trigger; the entire
pairing/IBE/threshold/envelope layer is reused unchanged (§5,
{[}drand-analysis §5, §8{]}).

drand itself is run by the \textbf{League of Entropy}, a federation of
more than a dozen independent organizations (founded 2019 by Protocol
Labs, Cloudflare, EPFL, Kudelski, and the University of Chile, among
others; membership has since rotated) --- the operational template for
Warden's public-good federation (§9).

\subsection{2.3 Time-lock puzzles and
VDFs}\label{time-lock-puzzles-and-vdfs}

A different lineage realizes delay through \emph{sequential
computation}: time-lock puzzles {[}Rivest--Shamir--Wagner 1996{]} and
verifiable delay functions {[}Wesolowski 2019; Pietrzak 2019{]} force a
solver to perform an inherently sequential amount of work before a value
is recoverable. These are \emph{self-contained} (no committee) but pay
for it: the ``time'' is wall-clock compute that must actually be
expended, the horizon is hard to calibrate across hardware generations,
and --- decisively for our setting --- they are \textbf{temporal, not
conditional}, and assume a present solver. They do not fit a
decades-dormant, owner-absent, condition-gated release. Warden is a
committee-release design, not a sequential-work one; we note the
contrast rather than compete on it.

\subsection{2.4 Witness encryption --- the trustless ideal, and why not
yet}\label{witness-encryption-the-trustless-ideal-and-why-not-yet}

\textbf{Witness encryption} {[}Garg--Gentry--Sahai--Waters 2013{]}
encrypts to an NP statement such that anyone holding a \emph{witness}
can decrypt and no one else can. Applied to our setting, the on-chain
proof that \texttt{executed(beatId)==true} would \emph{itself} be the
decryption key --- \textbf{no federation at all}, hence genuinely
trustless timing. This is the endgame. It is not deployable: general WE
rests on indistinguishability obfuscation or similarly idealized
assumptions, and even recent \textbf{WE-from-SNARKs/KZG} results
{[}e.g.~ePrint 2025/1364{]} that target concrete blockchain state remain
theoretical (they require Base's \texttt{executed} flag expressed as a
SNARK/KZG statement under idealized assumptions). Warden treats WE as
the gen-3 successor (§10) and keeps its envelope \texttt{alg}-versioned
so a WE outer layer can drop in with no change to recipients or to
Maktub.

\textbf{Signature-based witness encryption (SWE)} {[}McFly, ePrint
2022/433; compact SWE 2024/1477; eWEB 2025/1064{]} is a nearer
approximation: encrypt against a committee's verification keys plus a
\emph{tag}, such that the committee's threshold \emph{signature on the
tag} is the decryption witness --- users never interact with the
committee, and \emph{no per-item decryption key and no master secret are
ever held.} This ``no master secret'' property is, on paper, the
longevity-optimal design for a dead-man's-switch, and it runs on the
same BLS12-381 substrate. Warden targets SWE for gen-2 (§10) but does
not build on it today: the only implementations are dormant research
prototypes, ciphertext is linear in committee size (heavy on mobile),
and the compact variant requires iO.

\subsection{2.5 Removing DKG, and detecting
collusion}\label{removing-dkg-and-detecting-collusion}

Two recent advances bear directly on Warden's two worst liabilities (DKG
complexity; undetectable collusion). \textbf{Silent threshold
encryption} {[}ePrint 2024/263{]} derives the joint public key as a
deterministic function of operators' \emph{locally published} keys ---
no interactive DKG ceremony. \textbf{Silent threshold traitor-tracing}
(ST3) {[}ePrint 2025/342; CCS 2025{]} adds \emph{public tracing of a
colluding committee with no trusted authority}. The latter is
significant for us: gen-1 Warden cannot detect silent early-reveal
collusion (weakness W1, §6), and traitor-tracing would make such
collusion \emph{publicly attributable}, converting an undetectable
attack into a detectable, deterrable one. We treat traitor-tracing as a
\textbf{layer} to adopt when it has audited code, independent of the
base scheme (§8 of the architecture decision).

\subsection{2.6 Deployed and adjacent conditional-decryption
systems}\label{deployed-and-adjacent-conditional-decryption-systems}

Warden is not the first system to release decryption capability on a
condition; positioning it honestly against deployed and academic
neighbours sharpens the actual delta.

\begin{itemize}
\tightlist
\item
  \textbf{TACo (Threshold Access Control), Threshold Network / NuCypher}
  --- a productized threshold network that releases decryption rights
  when access \emph{conditions}, including on-chain state, are
  satisfied. This is the closest \emph{productized} competitor (now
  under World Ethical Data Foundation stewardship, with a stable
  relaunch slated for 2026); its conditions are point-in-time access
  checks rather than latching triggers. Warden's distinguishing
  requirements are the \emph{monotone-latching} condition discipline
  (§3.1), \emph{permanence across disjoint operator sets via resharing}
  over a decades horizon (§5.4), and an explicit \emph{finality policy}
  for the chain read (§4.4).
\item
  \textbf{Shutter Network} --- threshold encryption with ``keypers''
  releasing decryption keys on a trigger (originally anti-front-running
  for mempools). It shares the threshold-release shape but targets
  short-horizon, per-epoch keys, not dormant, owner-absent,
  resettable-condition ciphertexts.
\item
  \textbf{Calypso} {[}Kokoris-Kogias et al.~2021{]} --- on-chain secret
  management that releases threshold-encrypted secrets when an on-chain
  access policy is met. Direct academic prior art for ``decrypt when a
  blockchain condition holds''; Warden differs in the owner-absent
  dead-man's-switch trigger, the decades-permanence requirement, and the
  multi-platform client.
\item
  \textbf{Encryption to the Future / YOSO} {[}Campanelli et al.~2022;
  Gentry et al.{]} --- sending ciphertexts to \emph{future} committees
  selected over time; it motivates the owner-absent, operator-churn
  setting that resharing addresses here.
\end{itemize}

The point is not that conditional decryption is new --- it is not ---
but that the specific \emph{combination} in §2.7 is not served by these
systems. \emph{(The precise competitive deltas, especially against
TACo's evolving condition language, should be re-verified against each
system's current capabilities before any future peer-reviewed version.)}

\subsection{2.7 How Warden differs}\label{how-warden-differs}

Against all of the above, Warden's distinguishing combination is:
\textbf{(a)} a \emph{resettable on-chain condition} as the release
trigger (vs.~fixed round/time in tlock and time-lock puzzles);
\textbf{(b)} an explicit \emph{decades-horizon, owner-absent} threat
model with permanence-through-resharing; and \textbf{(c)} a deployed
multi-platform system (node + WASM + mobile FFI) with a finality policy
for the chain read --- a security-critical surface that has no analog in
drand because drand has a clock, not a chain to observe. None of these
is new \emph{cryptography}; together they are a new \emph{system}.

\begin{center}\rule{0.5\linewidth}{0.5pt}\end{center}

\section{3. Problem Statement and Threat
Model}\label{problem-statement-and-threat-model}

\subsection{3.1 The resettable-condition
trigger}\label{the-resettable-condition-trigger}

Let \texttt{C} be a \textbf{condition}: a canonical, public,
deterministic predicate over finalized on-chain state. The release
predicate Warden enforces is:

\begin{quote}
\emph{Release the decryption key for ciphertext \texttt{e} if and only
if \texttt{C(e)} is observed} true \emph{at finalized block depth.}
\end{quote}

Two properties make \texttt{C} suitable, both load-bearing (§4.2):

\begin{itemize}
\tightlist
\item
  \textbf{Determinism.} Every honest operator, reading finalized chain
  state at the agreed confirmation depth, must compute the \emph{same}
  yes/no. A threshold scheme cannot release if honest parties disagree
  on whether the condition holds.
\item
  \textbf{Monotonicity (latching).} Once \texttt{C} becomes true it
  stays true. Because a released key cannot be un-released, the trigger
  must not be allowed to flip back. The canonical instance ---
  \texttt{executed(beatId)==true} for a Maktub Beat --- is naturally
  latching: the contract never un-executes. Non-monotone raw predicates
  (e.g.~\texttt{price\ \textgreater{}\ X}) are rejected unless
  \emph{anchored} into a permanent form (``became true at/by block B'',
  or an event-was-emitted form).
\end{itemize}

The \emph{resettable} character lives entirely on the \textbf{sender's}
side and entirely \emph{before} the trigger: the owner keeps \texttt{C}
false by checking in, and revocation is the owner constructing
\texttt{C} so that it can be made \textbf{permanently unsatisfiable}
(Maktub's \texttt{deactivate(beatId)} ensures \texttt{executed} can
never become true → Warden never releases → the ciphertext is permanent
gibberish). From Warden's perspective the \emph{observed} predicate is
still monotone; resettability is a property of how the application
drives the underlying contract state.

\subsection{3.2 Actors}\label{actors}

\begin{itemize}
\tightlist
\item
  \textbf{Author / owner} --- encrypts the payload offline and drives
  the condition (checks in, or revokes). Absent at release time by
  assumption.
\item
  \textbf{Recipient(s)} --- hold the private key that the inner layer is
  encrypted to; the only parties who can ever read content.
\item
  \textbf{Operator federation} --- \texttt{n} independent organizations,
  each holding one share of the master key; each runs a node that
  observes the chain and serves partial decryptions.
\item
  \textbf{Chain} --- the public ledger (Base) whose finalized state
  defines the condition.
\item
  \textbf{Adversary} --- see §3.4.
\end{itemize}

\subsection{3.3 Security goals}\label{security-goals}

\begin{itemize}
\tightlist
\item
  \textbf{G1 --- Time-confidentiality.} Before \texttt{C} holds, the
  payload is unreadable by \emph{everyone, including the intended
  recipient.}
\item
  \textbf{G2 --- Delivery availability.} After \texttt{C} holds, the
  payload is readable by the intended recipient, and remains so
  indefinitely (release is idempotent and retryable forever). In the
  current client-requested mode (§5.3), release is
  \emph{recipient-initiated} --- a recipient must pull and combine ---
  so this is availability-on-demand, not push delivery; the autonomous
  mode (§5.3) pre-produces partials node-side, but a recipient still
  fetches and decrypts.
\item
  \textbf{G3 --- Content confidentiality.} Only the intended recipient
  ever reads plaintext --- \emph{including} against a fully-colluding
  federation.
\end{itemize}

\subsection{3.4 Adversary model and
assumptions}\label{adversary-model-and-assumptions}

We assume:

\begin{itemize}
\tightlist
\item
  \textbf{Threshold-bounded corruption} for timing.
  Confidentiality-before-trigger (G1) holds as long as fewer than
  \texttt{t} operators collude; availability-after-trigger (G2) holds as
  long as at least \texttt{t} operators are honest and live. With
  \texttt{t\ ≈\ ⌈2n/3⌉} these bounds trade off (raising \texttt{t}
  strengthens G1 but weakens G2-liveness, and vice versa) --- the
  standard threshold tension, stated precisely in §6.1; it is \emph{not}
  an ``honest majority''. \textbf{G3 rests on neither} (it is pure
  mathematics, §6).
\item
  \textbf{Correct finalized-state observation:} honest operators read
  true finalized chain state (not an eclipsed or spoofed view). This is
  part of the trust model and is \emph{new relative to drand}, which has
  no chain to observe (§4.4, §6/W3, §6/W4).
\item
  \textbf{Recipient key secrecy:} the recipient's private key is not
  compromised (otherwise G3 is moot --- but only for that one message).
\item
  \textbf{Independence of operators:} the security parameter is not raw
  node count but \emph{genuine independence} (distinct organizations,
  jurisdictions, and infrastructure), so that no single subpoena,
  coercion, or failure reaches \texttt{t} (§6/W4, §9).
\end{itemize}

The two-axis trust split is the heart of the model and is stated in full
in §6.1.

\begin{center}\rule{0.5\linewidth}{0.5pt}\end{center}

\section{4. System Design}\label{system-design}

\subsection{4.1 The double-wrap
envelope}\label{the-double-wrap-envelope}

Warden's defining design choice is that the federation's power is
\emph{confined to timing} and can never extend to content. This is
achieved by nesting two independent gates over a single random content
key \texttt{K}:

\begin{verbatim}
K          = random symmetric content key
inner      = AEAD(payload, K)                       ← the content
K_wrapped  = ECIES(K, recipientPub)                 ← recipient-gated.  PURE MATH.
outer      = Warden-IBE(K_wrapped, H(condition))    ← condition-gated.  Federation can open this.
\end{verbatim}

Warden's release only ever opens \textbf{outer}, which reveals
\texttt{K\_wrapped} --- \emph{an ECIES ciphertext}, not \texttt{K} and
not the payload. Only the recipient's private key turns
\texttt{K\_wrapped} into \texttt{K}. Therefore \textbf{both} the
condition (to open \texttt{outer}) \textbf{and} the recipient key (to
open \texttt{K\_wrapped}) are required, and \textbf{even a
fully-colluding federation that releases every outer key learns only
release timing, never content} (this is goal G3). Content
confidentiality is independent of Warden entirely (given recipient-key
secrecy, §3.4).

\textbf{As implemented.} Because the Boneh--Franklin IBE plaintext block
is a fixed 32 bytes, the outer layer is realized as a hybrid:
\texttt{outer.ibe\ =\ IBE(obk,\ H(condition))} gates a random 32-byte
\texttt{obk}, and \texttt{outer.seal\ =\ AEAD\_obk(K\_wrapped)} seals
the ECIES-wrapped key under \texttt{obk}. This is logically identical to
``IBE-gate \texttt{K\_wrapped}''. Hardening from a security review: both
AEAD layers bind \texttt{domain\ ‖\ network\ ‖\ H(condition)} as
associated data (tamper-evident and domain-separated), and the payload
is \textbf{bucket-padded} to hide its length. AEAD is ChaCha20-Poly1305
{[}RFC 8439{]}; the recipient layer is ECIES over secp256k1 --- chosen
so recipients reuse their existing Ethereum wallet keypair rather than
manage a second key --- with HKDF info bound to
\texttt{ephPub\ ‖\ recipientPub} (SEC1 shared-info). (The IBE/threshold
layer is on BLS12-381; the two curves serve the two different layers.)
\textbf{No recipient metadata is stored in the envelope} --- the
recipient is implicit in who can ECIES-open \texttt{K\_wrapped}. All
domain tags and pad buckets are provisional and will be frozen with
cross-language test vectors before mainnet.

The envelope is \texttt{alg}-versioned (\texttt{warden-v1}); future
schemes (SWE, witness encryption) replace only the
\texttt{outer}/\texttt{condition} handling, leaving \texttt{inner} and
the integration unchanged, and old envelopes remain decryptable forever
(§10). An optional age-stanza form
(\texttt{-\textgreater{}\ warden\ \textless{}network-hash\textgreater{}\ \textless{}H(condition)\textgreater{}})
provides interoperability with the tlock/age ecosystem, mirroring
tlock's stanza with \emph{round} → \emph{H(condition)} and
\emph{chainHash} → Warden network hash.

\subsection{4.2 The condition model}\label{the-condition-model}

The IBE identity a payload is encrypted to is the hash of a
\textbf{canonical, versioned condition specification}, serialized with
\textbf{RFC 8785 (JSON Canonicalization Scheme)} so the bytes are
byte-identical across the Rust, TypeScript, and Dart implementations:

\begin{verbatim}
identity = H( "warden-cond-v1" ‖ jcs_rfc8785(condition) )
\end{verbatim}

When asked for its partial on \texttt{identity}, a node is given the
condition \texttt{C} (public metadata in the envelope), checks
\texttt{H(C)\ ==\ identity}, independently evaluates \texttt{C} against
finalized chain state, and returns its partial \emph{iff} \texttt{C}
holds. This \textbf{binds the condition cryptographically to the
ciphertext}: conditions cannot be swapped after encryption. The
condition is \emph{public} --- it is metadata, not a secret (for Veil it
reveals \texttt{beatId} and that release gates on the beat's status,
both already public on-chain --- no new leakage).

The schema (\texttt{cond-v1}) supports exactly one \texttt{type} per
condition plus a mandatory \texttt{meta}:

\begin{itemize}
\tightlist
\item
  \texttt{contract} ---
  \texttt{\{\ chain,\ address,\ fn,\ args,\ word?,\ test,\ meta\ \}}
\item
  \texttt{block} ---
  \texttt{\{\ chain,\ field:\ "number"\textbar{}"timestamp",\ cmp,\ value,\ meta\ \}}
\item
  \texttt{event} ---
  \texttt{\{\ chain,\ address,\ sig,\ args,\ meta\ \}}
\item
  \texttt{all\ \textbar{}\ any\ \textbar{}\ not\ \textbar{}\ threshold}
  ---
  \texttt{\{\ of:\ {[}...sub-conditions{]},\ k\ (threshold\ only),\ meta\ \}}
\end{itemize}

The Veil release condition for a real Maktub Beat illustrates a subtlety
that the canonicalization must survive: \texttt{MaktubCore} exposes
\textbf{no \texttt{executed(uint256)} getter} --- execution status is
field 7 (the 8th element) of the \texttt{getHeartbeat(uint256)} return
tuple --- so the condition selects it with \texttt{word:\ 7}:

\begin{Shaded}
\begin{Highlighting}[]
\FunctionTok{\{}
  \DataTypeTok{"type"}\FunctionTok{:} \StringTok{"contract"}\FunctionTok{,} \DataTypeTok{"chain"}\FunctionTok{:} \DecValTok{84532}\FunctionTok{,}
  \DataTypeTok{"address"}\FunctionTok{:} \StringTok{"0xb603C96D089F64Ac487EE0bdaE97D49848F86133"}\FunctionTok{,}
  \DataTypeTok{"fn"}\FunctionTok{:} \StringTok{"getHeartbeat(uint256)"}\FunctionTok{,} \DataTypeTok{"args"}\FunctionTok{:} \OtherTok{[}\StringTok{"123"}\OtherTok{]}\FunctionTok{,} \DataTypeTok{"word"}\FunctionTok{:} \DecValTok{7}\FunctionTok{,}
  \DataTypeTok{"test"}\FunctionTok{:} \FunctionTok{\{} \DataTypeTok{"cmp"}\FunctionTok{:} \StringTok{"=="}\FunctionTok{,} \DataTypeTok{"value"}\FunctionTok{:} \KeywordTok{true} \FunctionTok{\},}
  \DataTypeTok{"meta"}\FunctionTok{:} \FunctionTok{\{} \DataTypeTok{"finality"}\FunctionTok{:} \DecValTok{32}\FunctionTok{,} \DataTypeTok{"tier"}\FunctionTok{:} \DecValTok{1} \FunctionTok{\}}
\FunctionTok{\}}
\end{Highlighting}
\end{Shaded}

Two serialization rules are load-bearing because the bytes are hashed
into the identity, and a mismatch yields a \emph{different} identity
that silently breaks decryption: (i) \textbf{all \texttt{uint256}
args/values are decimal strings}, never JSON numbers, to avoid
JavaScript's \(2^{53}-1\) precision loss and to guarantee byte-identical
serialization across platforms; (ii) \texttt{word} is \textbf{omitted
from the canonical form when 0}, so single-value-getter conditions hash
identically to the pre-\texttt{word} schema.

\textbf{Trust tiers.} \emph{Tier 1} --- on-chain state
(\texttt{contract}/\texttt{block}/\texttt{event}, finalized) --- is
deterministic and is what v1 ships. \emph{Tier 2} --- external/oracle
data --- adds a data-source trust and a determinism risk; it is opt-in
and later-phase, and operators \emph{declare} which tiers they evaluate
so clients can decide (§6/W7). The same condition machinery generalizes
well beyond Veil --- time capsules
(\texttt{block.timestamp\ \textgreater{}=\ T}), multi-trigger delivery
(\texttt{any(executed,\ deadline)}), milestone escrow, DAO-vote-gated
reveals, or any third-party EVM contract getter --- which is the
adoption story (§9).

\subsection{4.3 Network planes}\label{network-planes}

Warden inherits drand's hard-won plane separation. drand's v2.0
post-mortem documented a reshare-downgrade attack that arose precisely
from \emph{merging} operator-control RPC with node-to-node RPC on one
service; the fix split them, and Warden adopts the separation from the
start:

\begin{itemize}
\tightlist
\item
  \textbf{Control plane} --- local-only operator management
  (\texttt{wardenctl} → daemon over a localhost port). \textbf{Never
  internet-exposed.}
\item
  \textbf{Private plane} --- node-to-node messaging for DKG and
  resharing, authenticated and behind TLS (terminated by a reverse
  proxy; the daemon does not terminate TLS itself).
\item
  \textbf{Public plane} --- clients fetch partial decryptions over
  HTTP/relay. Read-only; partials are publicly verifiable against the
  master public key, so the relay need not be trusted.
\end{itemize}

\subsection{4.4 Finality policy}\label{finality-policy}

A node must decide \emph{when} an on-chain condition is ``final enough''
to release on. This is the genuinely new, security-critical surface
Warden adds over drand, because \textbf{a reorg that reverts the
condition after a release is an irreversible confidentiality leak}
(§6/W3). Policy:

\begin{itemize}
\tightlist
\item
  Read at the chain's \textbf{finalized/safe head}, or at least the
  chain's finality depth (for Base, prefer L1-finalized state). The
  current implementation reads at the \texttt{finalized} tag.
\item
  The \textbf{minimum finality floor is a federation-wide consensus
  parameter --- identical for every node, not configured locally.}
  Per-node floors would make nodes disagree on whether a condition is
  final and stall release (a node with a higher local floor would block
  the threshold). \texttt{C.meta.finality} may override the floor
  \emph{upward only}, never below it.
\item
  The residual reorg risk for the chosen depth is documented, not
  hidden.
\end{itemize}

\begin{center}\rule{0.5\linewidth}{0.5pt}\end{center}

\section{5. Cryptographic
Construction}\label{cryptographic-construction}

The cryptography is reused from drand near-verbatim; what follows states
it precisely and marks the one new ingredient (the condition gate).

\subsection{5.1 Setup (key generation)}\label{setup-key-generation}

\begin{itemize}
\tightlist
\item
  \textbf{Mainnet --- real DKG, no dealer.} All operators run one DKG
  ceremony (Pedersen/Feldman VSS, drand-derived), producing one
  \textbf{master public key} \texttt{P\_pub} and one \textbf{share
  \texttt{sk\_i} per operator}; the master private key \texttt{msk} is
  never assembled. The threshold targets a BFT-style
  \texttt{t\ ≈\ ⌈2n/3⌉} (minimum \texttt{⌊n/2⌋+1}). The deployment
  inherits drand's lifecycle (joiners/remainers/leavers), strict
  control-/node-plane separation, and \textbf{\texttt{OldThreshold}
  downgrade protection} on every reshare (the v2.0 attack class).
\item
  \textbf{Testnet --- trusted dealer.} A one-off script generates
  \texttt{msk}, Shamir-splits it into \texttt{n} proper field-element
  shares, and injects one per node. The full key briefly exists in the
  dealer --- \textbf{testnet-only, never mainnet} --- and the script
  outputs \texttt{P\_pub} (client config) and optionally each node's
  share-public-key (for partial verification).
\end{itemize}

\subsection{5.2 Encryption (client, offline, no network
interaction)}\label{encryption-client-offline-no-network-interaction}

\begin{verbatim}
K          = random symmetric content key
inner      = AEAD(payload, K)
K_wrapped  = ECIES(K, recipientPub)
identity   = H("warden-cond-v1" ‖ jcs_rfc8785(condition))
outer      = IBE_encrypt(K_wrapped, identity, P_pub)      // BF-IBE (U, V, W)
envelope   = { alg: "warden-v1", network, condition, outer, inner }
\end{verbatim}

The Boneh--Franklin ciphertext (with the FO transform for CCA security)
is, following tlock:

\begin{verbatim}
U = r·G
V = σ  XOR  H2( e(Q_id, P_pub)^r )      where Q_id = H1(identity)
W = M  XOR  H4(σ)                        with r = H3(σ ‖ M)   (Fujisaki–Okamoto)
\end{verbatim}

The pairing is written for the concrete BLS12-381 placement used by
tlock/drand quicknet: the identity hashes into G1
(\texttt{Q\_id\ =\ H1(identity)\ ∈\ G1}), the master public key lives in
G2 (\texttt{P\_pub\ ∈\ G2}), and partials/decryption keys are in G1,
with \texttt{e:\ G1\ ×\ G2\ →\ G\_T}. Because the BF plaintext block is
a fixed 32 bytes, \texttt{outer} is realized as the hybrid of §4.1 ---
\texttt{IBE(obk)} gates a random 32-byte key that AEAD-seals the larger
\texttt{K\_wrapped} --- which is logically identical to IBE-gating
\texttt{K\_wrapped} directly, but is \emph{necessary} rather than
cosmetic since \texttt{K\_wrapped} exceeds one BF block.

\subsection{5.3 Release (the condition gate --- the new
component)}\label{release-the-condition-gate-the-new-component}

Each node, on a request for \texttt{identity} carrying condition
\texttt{C}:

\begin{enumerate}
\def\labelenumi{\arabic{enumi}.}
\tightlist
\item
  Verify \texttt{H(C)\ ==\ identity} (binding).
\item
  Verify \texttt{C.type}/\texttt{tier} is supported by this node.
\item
  \textbf{Evaluate \texttt{C} against finalized chain state} (read the
  named chain at no fewer than \texttt{C.meta.finality} confirmations,
  never below the federation floor).
\item
  If \texttt{C} holds, return the partial
  \texttt{sig\_i\ =\ sk\_i\ ·\ H1(identity)} (publicly verifiable
  against node \texttt{i}'s share-public-key); otherwise return
  \texttt{ConditionNotMet}.
\end{enumerate}

The client collects \texttt{t} valid partials, Lagrange-combines them
\emph{in the exponent} into the IBE decryption key
\texttt{outerKey\ =\ msk\ ·\ H1(identity)}, IBE-decrypts \texttt{outer}
to recover \texttt{K\_wrapped}, has the recipient ECIES-open
\texttt{K\_wrapped} to \texttt{K}, and AEAD-decrypts \texttt{inner} to
the plaintext. Decryption (given a met condition) requires no further
interaction beyond gathering partials.

\textbf{Identity domain-separation (security-critical).} Unlike drand
round numbers, which are globally unique by construction,
condition-derived identities are \emph{attacker-influenceable}. To
prevent cross-item key reuse and grinding, the hash-to-curve \texttt{H1}
must be strongly domain-separated and must bind the chain id, contract
address, and full condition. A partial minted for one condition must
never open a ciphertext for another (§6/W7).

\textbf{Idempotent and retryable.} Because the condition is monotone and
chain state is permanent, release can be requested at any later time and
always succeeds once met; nodes may cache released partials and serve
them forever. A temporarily-offline network therefore causes
\emph{delay, not loss} (§6/W2). The current implementation ships
\textbf{client-requested} signing (\texttt{POST\ /partial} evaluates and
signs on demand); an \textbf{autonomous} mode (nodes watch the contract
and pre-produce partials when the condition flips, for better liveness)
is the mainnet target.

\subsection{5.4 Resharing --- the permanence
mechanism}\label{resharing-the-permanence-mechanism}

When operators join or leave (or for periodic share refresh), the
federation runs \textbf{resharing}: each old node deals its existing
share as the \emph{free coefficient} of a new polynomial; Lagrange
interpolation gives the new node set shares of the \textbf{same}
\texttt{msk} --- so \texttt{P\_pub} is unchanged and \textbf{all
historical ciphertexts remain decryptable}. Old and new operator sets
may be entirely disjoint. \texttt{OldThreshold} MUST be specified
(downgrade protection). Resharing \emph{can} change membership, every
share (rotation), and the threshold; it \emph{cannot} change
\texttt{P\_pub} (and hence the decryptability of past ciphertexts). This
is precisely the \textbf{permanence-with-churn} property that no
off-the-shelf vendor offered and that the decades-horizon requirement
demands.

\subsection{5.5 Verifiability}\label{verifiability}

Every partial \texttt{sig\_i} is a BLS signature verifiable against node
\texttt{i}'s published share-public-key; the combined key verifies
against \texttt{P\_pub}. Clients reject invalid partials before
combining (relay-trust-free, exactly as drand clients verify beacons
against the chain public key). A partial released for a condition that
is \emph{not} met is therefore \textbf{publicly provable misbehavior},
which can feed an eviction/slashing process in tokenized deployments
(not required for the public-good federation, §9).

\begin{center}\rule{0.5\linewidth}{0.5pt}\end{center}

\section{6. Security Analysis}\label{security-analysis}

This section is the credibility core. We state the trust split exactly
and then walk the weaknesses one by one; for this audience, the rigor of
the limitation analysis \emph{is} the contribution's value.

\subsection{6.1 The two-axis trust
split}\label{the-two-axis-trust-split}

\begin{longtable}[]{@{}
  >{\raggedright\arraybackslash}p{(\linewidth - 4\tabcolsep) * \real{0.3333}}
  >{\raggedright\arraybackslash}p{(\linewidth - 4\tabcolsep) * \real{0.3333}}
  >{\raggedright\arraybackslash}p{(\linewidth - 4\tabcolsep) * \real{0.3333}}@{}}
\toprule\noalign{}
\begin{minipage}[b]{\linewidth}\raggedright
Property
\end{minipage} & \begin{minipage}[b]{\linewidth}\raggedright
Rests on
\end{minipage} & \begin{minipage}[b]{\linewidth}\raggedright
Strength
\end{minipage} \\
\midrule\noalign{}
\endhead
\bottomrule\noalign{}
\endlastfoot
\textbf{Content confidentiality} (only the recipient ever reads
plaintext --- G3) & the recipient's own key (ECIES \emph{inner} layer) &
\textbf{Pure math.} Even a fully-colluding federation cannot read
content. \\
\textbf{Timing + revocation} (unreadable until the condition holds ---
G1/G2) & a \textbf{threshold-bounded} federation (fewer than \texttt{t}
collude) \emph{and} nodes correctly observing the on-chain condition &
\textbf{Threshold}, not trustless. \\
\end{longtable}

The double-wrap is what confines the federation's power to
\emph{timing}: a colluding federation can release the outer key early,
or refuse it, but never read content.

The timing axis has two distinct thresholds, not one ``majority''.
Confidentiality (G1) is broken only if \texttt{t} operators collude to
release early; availability (G2) fails only if fewer than \texttt{t}
honest operators are reachable to release at all. With
\texttt{t\ ≈\ ⌈2n/3⌉} the design tolerates up to \texttt{t−1} colluders
without an early leak and up to \texttt{n−t} unavailable nodes without
stalling --- and the two cannot both be maximized, so the choice of
\texttt{t} is an explicit confidentiality-versus-liveness trade, not a
single honest-majority assumption.

\subsection{6.2 Weaknesses (extensive, for the deployable
design)}\label{weaknesses-extensive-for-the-deployable-design}

\textbf{W1 --- Early reveal by bribing a threshold.} A \emph{receiver}
who corrupts \texttt{t} operators can have the outer key released before
the condition holds and --- with their own recipient key --- read early.
Bounds and defenses: only the receiver benefits (content is
inner-locked, so a third-party briber gains nothing readable); the cost
is corrupting \texttt{t} \emph{independent} operators, so the defense is
economic and diversity-based (a large, jurisdictionally-diverse
federation and a higher \texttt{t}); the attack bites hardest on
high-value beats (large-wallet inheritance) where the bribe bar equals
the wallet's value. An optional trigger-gated-identity construction can
force such an attack to be a \emph{loud, detectable blanket release}
rather than a quiet targeted one. As of the 2026-06 scan,
\textbf{threshold traitor-tracing} (ST3, ePrint 2025/342) can make a
colluding committee \emph{publicly attributable} --- turning W1 from
undetectable to detectable-and-deterrable --- and is the tracked
mitigation (a layer to adopt when it has audited code, §10).

\textbf{W2 --- Liveness / denial → mostly delay, not loss.} Because the
ciphertext and the condition are permanent, release is retryable
forever; a temporarily down or refusing network causes \emph{delayed}
delivery, not loss, and the receiver can pull at any later time.
Permanent loss occurs only if \emph{all} redundant Warden
generations/networks die forever --- mitigated by resharing
(permanence-with-churn) plus multi-network redundancy; there is
deliberately no first-party backup key (see W4). Delay severity is
use-case dependent: acceptable for inheritance, harmful for urgent
safety triggers, which argues for high availability there.

\textbf{W3 --- Finality / reorg leak (new; no drand precedent).} If a
node releases on an \texttt{executed==true} state that is subsequently
reorged away, the reveal is irreversible. Mitigation is the finality
policy of §4.4 (release only on finalized state at a conservative,
federation-wide depth). This is a genuinely new, security-critical
surface that Warden adds over drand precisely because Warden observes a
chain rather than a clock.

\textbf{W4 --- Operator/infra capture, and the no-first-party-key rule.}
A federation entirely controlled by one party (e.g.~all nodes in one
cloud account) collapses to a single party: compellable and capable of
early release. \textbf{Warden never holds a release key itself}, and
real security requires \emph{genuinely independent} operators. This is
why the present testnet --- five nodes, but a single trust domain (one
operator, one provider) --- has \textbf{zero timing-security and is for
protocol exercise only.} A related capture vector is
\emph{condition-source integrity}: if many operators rely on the
\emph{same} public RPC provider, a compromise or \textbf{eclipse} of
that provider can spoof \texttt{executed==true} and trick a threshold
into early release \emph{without breaking any cryptography}. Mainnet
operators must therefore run their own chain node or use independent,
authenticated, jurisdictionally-diverse RPC endpoints, and must read at
the federation-wide finality floor. The drand v2.0 lesson is inherited
in full: strict control-/node-plane separation and \texttt{OldThreshold}
downgrade protection on every reshare, or a malicious operator can chain
reshares to reduce the threshold to themselves and recover \texttt{msk}.

\textbf{W5 --- Master-key permanence is double-edged.} \texttt{P\_pub}
is immutable across reshares (the permanence feature). The consequence:
\emph{any} eventual reconstruction of \texttt{msk} breaks \emph{all}
past and future ciphertexts under that key --- and Warden's ciphertexts
may sit dormant for years. Reshare hygiene is non-negotiable, and a
key-rotation/new-generation migration story (opt-in re-encryption,
mirroring Maktub's immutable-V2 pattern) is part of the roadmap.

\textbf{W6 --- Not quantum-resistant.} BLS/IBE are pairing-based, not
post-quantum (the stated tlock horizon is on the order of five years).
For payloads meant to ``survive you by decades'' this is a real limit to
disclose to product and legal stakeholders. There is currently no
deployable post-quantum replacement for threshold-IBE/BLS, so the
decades claim is unbacked \emph{at the threshold layer}; the mitigation
is to keep the envelope \texttt{alg}-versioned and to
post-quantum-harden the \emph{inner} recipient layer first, where
deployable PQ KEMs exist (§10).

\textbf{W7 --- Condition-evaluation trust (Tier-2).} Tier-1 (on-chain,
finalized) is deterministic. Tier-2 (oracle/API data) adds a data-source
trust and a determinism risk; it stays opt-in and clearly labeled, with
operators declaring supported tiers. Identity domain-separation (§5.3)
must hold so that a partial for one condition cannot open another
(grinding/replay).

\textbf{W8 --- Open-request metadata channel (client → node).} Opening a
Veil envelope POSTs the \emph{condition} (which embeds \texttt{beatId},
the core address, and the chain id) to every node's \texttt{/partial}.
Each node operator --- and any on-path observer --- therefore learns
\emph{the link between a requester IP and a beatId of interest} and the
\emph{timing} of the open attempt. The \texttt{beatId} and condition are
already public on-chain; what is new is binding them to a requesting IP
at open time. This is consistent with ``not metadata-private'' (§6.3),
not a content leak (a tampered condition simply never releases ---
fail-closed). Mitigations for the independent-operator federation:
client egress over a privacy-preserving transport (e.g.~Tor or an
oblivious relay) and TLS pinning / a known-key channel to deny passive
collectors the linkage. With the all-ours testnet federation this is
moot, but it must be addressed before independent operators run.

\textbf{W9 --- Cross-sender identity collision.} The identity is
\texttt{H("warden-cond-v1"\ ‖\ jcs(condition))} with nothing per-message
mixed in, so two independent senders who encrypt to the \emph{same}
condition share an identity and therefore share the released outer key
\texttt{msk\ ·\ H1(identity)}. For Veil this is benign: every beat
carries a unique \texttt{beatId} in the condition \texttt{args}, so
identities never collide. But for the generalized conditions of §4.2
(e.g.~a popular time-capsule date used verbatim by many senders), an
identical condition couples all its users --- one requester's pull
releases the outer key for every ciphertext under that condition.
Content stays safe (the inner ECIES layer is per-recipient), but the
\emph{timing gate} is shared. Mitigation: applications that reuse
generic conditions should bind a per-message salt into the condition
\texttt{meta} (a node still hashes and recomputes the identity
deterministically), or otherwise engineer identity uniqueness as Veil
does via \texttt{beatId}.

\textbf{W10 --- Share-key registration / rogue-key.} Public
verifiability (§5.5) rests on each node's published share-public-key.
For the ``publicly provable misbehavior'' property to hold,
share-public-keys MUST be bound to the DKG transcript (or, on testnet,
the dealer's commitment) and registered with proof-of-possession;
otherwise an operator publishing an adversarially-chosen
share-public-key could undermine partial verification, as in classic BLS
rogue-key attacks. This is a requirement on the setup ceremony, stated
here so it is not assumed away. (The current trusted-dealer testnet
build is not exposed to this: the dealer computes every share-public-key
honestly by construction; the requirement bites once nodes self-publish
keys under DKG.)

\textbf{W11 --- Request-channel denial-of-service.} Each
\texttt{/partial} request triggers a finalized chain read and a BLS
signing operation, so an unauthenticated public endpoint is a
compute/RPC amplification vector. This is an availability concern, not a
confidentiality one (an unmet or malformed condition simply yields no
release --- fail-closed). Mitigation: rate-limiting, lightweight request
authentication, and caching of already-released partials at the public
plane.

\subsection{6.3 What Warden is not trying to
be}\label{what-warden-is-not-trying-to-be}

\begin{itemize}
\tightlist
\item
  \textbf{Not trustless.} The trustless ideal is witness encryption (the
  on-chain proof \emph{is} the key, no federation); it is not deployable
  today. Warden is the best deployable approximation and is designed to
  migrate to WE via the \texttt{alg} envelope when it ships.
\item
  \textbf{Not metadata-private.} Existence, timing, condition, and (for
  Veil) \texttt{beatId}, recipient set, and the sender-recipient
  relationship are public. Warden hides \emph{content until the
  trigger}, never \emph{that something exists}.
\end{itemize}

\subsection{6.4 Honest claim language}\label{honest-claim-language}

For any derived copy: one \emph{may} say \emph{time-bound, revocable,
end-to-end-sealed delivery enforced by an independent threshold
federation --- secure against early release unless \texttt{t} operators
collude, and against denial as long as \texttt{t} stay honest and live
--- with redundancy.} One \emph{must not} say \emph{absolute},
\emph{mathematically guaranteed} (for timing/revocation),
\emph{bulletproof}, \emph{absolute anonymity}, \emph{prevents bribery},
\emph{cannot leak}, or \emph{honest-majority} (the bound is the
threshold \texttt{t}, not a majority).

\begin{center}\rule{0.5\linewidth}{0.5pt}\end{center}

\section{7. Implementation}\label{implementation}

Warden is implemented as a single \textbf{Rust core}
(\texttt{warden-core}) --- chosen for memory safety, mature
pairing-crypto crates, and the ability to drive one codebase to three
client targets --- plus the surrounding daemon, dealer, CLI, and
bindings. The reference implementation is complete and deployed on
testnet (§8).

\begin{longtable}[]{@{}
  >{\raggedright\arraybackslash}p{(\linewidth - 4\tabcolsep) * \real{0.3333}}
  >{\raggedright\arraybackslash}p{(\linewidth - 4\tabcolsep) * \real{0.3333}}
  >{\raggedright\arraybackslash}p{(\linewidth - 4\tabcolsep) * \real{0.3333}}@{}}
\toprule\noalign{}
\begin{minipage}[b]{\linewidth}\raggedright
Component
\end{minipage} & \begin{minipage}[b]{\linewidth}\raggedright
Crate / artifact
\end{minipage} & \begin{minipage}[b]{\linewidth}\raggedright
Role
\end{minipage} \\
\midrule\noalign{}
\endhead
\bottomrule\noalign{}
\endlastfoot
Crypto core & \texttt{warden-core} & Threshold BF-IBE encrypt/decrypt,
threshold-BLS partial/combine, the \texttt{warden-v1} double-wrap
envelope, the federation file format. On \texttt{arkworks} BLS12-381.
Library only. \\
Dealer & \texttt{warden-dealer} & Trusted-dealer ceremony CLI:
materialize \texttt{msk}, Shamir-split, write \texttt{federation.json}
(public) + per-node share files (secret, mode 0600). \textbf{Testnet
only.} \\
Node & \texttt{warden-node} (\texttt{wardend}) & Condition-watcher +
\texttt{POST\ /partial} threshold release. Reads Base Sepolia at the
\texttt{finalized} tag. The security-critical evaluator is
\texttt{node/src/eval.rs}. \\
Client CLI & \texttt{warden-cli} (\texttt{warden}) & \texttt{keygen} /
\texttt{encrypt} (double-wrap → CID store) / \texttt{decrypt} (poll
federation → combine → open, retry-until-released). \\
WASM & \texttt{warden-wasm} & \texttt{wasm-bindgen} bindings over core
for the TypeScript SDK: \texttt{condition\_identity} /
\texttt{seal\_gated} / \texttt{open\_gated} / \texttt{combine}. \\
Mobile FFI & \texttt{warden-ffi} & Thin C-ABI over core for the Flutter
app via \texttt{dart:ffi}. Panics are caught at the boundary and
returned as \texttt{\{ok:false,…\}} JSON rather than aborting the host
app. Cross-compiles to an iOS xcframework and 4-ABI Android
\texttt{jniLibs}. \\
\end{longtable}

A reproducible 3-node federation is provided via
\texttt{docker-compose}. The cryptographic loop is verified by the
crate's own offline integration tests and by a local devnet end-to-end
harness, and the full loop has been run end-to-end on live Base Sepolia (§8).
The same core, compiled to native, WASM, and mobile FFI, is what closes
the multi-platform-client gap (notably Flutter mobile) that disqualified
off-the-shelf alternatives. Both the specification and the
implementation are public and MIT-licensed (see \emph{Artifact
availability} below); the audit gate (§9, §10) governs operator
recruitment and mainnet, not the openness of the source.

The reuse/replace boundary is clean: the entire
IBE/threshold-BLS/DKG/resharing/age-envelope stack is reused
near-verbatim from the drand lineage, while the genuinely new
engineering is the round-scheduler-to-Base-condition-watcher swap, the
domain-separated \texttt{H(condition)} identity, the finality handling
for the chain read, and the on-demand (vs.~periodic-firehose)
condition-keyed release with an idempotent released-key store.

\textbf{Artifact availability.} Warden is developed in the open as a
standalone primitive at \textbf{https://github.com/bytesbrains/warden}
(MIT), with an overview and operator information at
\textbf{https://warden.bytesbrains.com}. The specification
(\texttt{docs/}) is public; the implementation and reproducible
federation are in the repository.

\begin{center}\rule{0.5\linewidth}{0.5pt}\end{center}

\section{8. Deployment Experience /
Evaluation}\label{deployment-experience-evaluation}

Warden is deployed as a \textbf{five-node federation on fly.io} and runs
the full loop end-to-end against \textbf{Base Sepolia}, integrated with
the live Maktub protocol and a mobile client (in testing): create a beat
(with a check-in interval), seal a payload to the condition
\texttt{getHeartbeat(beatId).word{[}7{]}\ ==\ true}, let the timer lapse
so an executor flips the beat to executed, wait for the configured
confirmation depth, collect partials from the federation, combine, and
decrypt. In this configuration the system operates as expected. The
complementary negative path --- \texttt{deactivate(beatId)} → the
condition can never be satisfied → the payload is never released
(permanent gibberish) --- also holds. A local Hardhat devnet harness
(\texttt{evm\_increaseTime} to skip the timer) drives the same flow
without funds and is the ``Warden works for all conditions'' gate.

Crucially, this deployment exercises the \textbf{protocol, its
integration, and liveness --- not the timing-security trust model.} All
five nodes are operated by a single party (BytesBrains) on one provider
(fly.io), i.e.~a single trust domain; per W4 (§6) that collapses to one
party and has, by construction, no early-release resistance. The missing
ingredient is \emph{independence} --- distinct organizations,
jurisdictions, and infrastructure --- which is being actively recruited
(§9). A revised version will add on-chain transaction evidence and a
measurement pass.

We characterize this as \textbf{deployment experience} rather than a
full performance evaluation, which is appropriate (and, at
applied/systems venues, acceptable) for a pre-audit testnet system. The
compute profile is favorable by construction: a single threshold-BLS
partial is microseconds of work per request, and per-node resource needs
are modest (low single-digit vCPU, a few GB of RAM, tens of GB of disk
for the share, config, released-partial cache, and logs, over commodity
always-on infrastructure with a Base RPC). A rigorous measurement pass
--- partial-release latency, combine latency, the finality-wait
distribution, per-node footprint, and per-operator bandwidth --- is
future work and is the natural content for a future revision.

\begin{center}\rule{0.5\linewidth}{0.5pt}\end{center}

\section{9. Federation Governance and
Sustainability}\label{federation-governance-and-sustainability}

Warden is intended to run as \textbf{token-free public-good
infrastructure}, in the drand / League-of-Entropy tradition. Warden, and
the Maktub protocol it serves, are public-good initiatives of
BytesBrains Pte Ltd; Maktub is an open, unregistered on-chain protocol,
and the operator \emph{federation} is an open collaboration of
independent organizations rather than a single legal entity ---
``foundation''/``federation'' here name that collaboration, not an
incorporated body. That this is \emph{not} one entity is the point:
independence across organizations is exactly the security property (§6,
W4). The governance posture is deliberate and is itself a security
property:

\begin{itemize}
\tightlist
\item
  \textbf{No token; operators credited, not paid.} The absence of a
  payment/staking economy keeps operators \emph{disinterested} ---
  neutrality and the lack of a financial stake are exactly what make a
  federation hard to bribe or capture, which is the property W1/W4
  depend on.
\item
  \textbf{Institutional, not personal, commitment.} The real ask of an
  operator is not hardware but \emph{reliability, key-share custody, and
  longevity} --- running for years and participating in occasional
  resharing ceremonies. This is why \textbf{institutions that persist}
  (universities especially) are the ideal operators; the effort once set
  up is a few hours per month of part-time sysadmin attention, with no
  dedicated team and no special hardware.
\item
  \textbf{Independence is the security, not node count.} Even a small
  federation of genuinely independent operators (different
  organizations, jurisdictions, infrastructure) gives a meaningful
  \texttt{t}-of-\texttt{n}; the goal is diversity and longevity, not a
  data center. A federation that converges on one cloud account or one
  RPC provider has the security of one party (W4).
\item
  \textbf{Funding line.} Audits, permanence, and continued development
  are funded through a treasury and grants rather than protocol fees;
  the honest status (testnet, pre-audit) is disclosed throughout. The
  decade-scale funding question --- who pays for audits and continuity
  across the full dormancy horizon --- is openly unresolved, and the
  paper does not print durability claims it cannot back.
\end{itemize}

\textbf{Recruiting independence (an open invitation).} This independence
is being actively pursued. The current five-node testnet runs within a
single trust domain (BytesBrains, on one provider), and discussions are
underway with universities and other long-lived institutions to take
over genuinely independent shares. We extend this as an \textbf{open
invitation}: any institution able to commit to reliability, key-share
custody, and longevity is welcome to operate a founding node. That
diversity --- distinct organizations, jurisdictions, and infrastructure
--- not node count, is what gives Warden its timing-security (§6, W4),
and turning this paper into that invitation is a deliberate part of why
it is published (§11).

The federation's behavior has \textbf{no governance surface over its own
cryptography}: membership churns via resharing, but the master public
key, and hence the decryptability of every dormant ciphertext, is
outside anyone's power to change. This mirrors the immutability
invariant of the Maktub protocol layer it serves.

\begin{center}\rule{0.5\linewidth}{0.5pt}\end{center}

\section{10. Limitations and Future
Work}\label{limitations-and-future-work}

The weaknesses of §6 are the honest limitation set; the forward path
addresses them through the \texttt{alg}-versioned envelope, which lets
the cryptography migrate with \textbf{no change to recipients or to
Maktub} and with old envelopes decryptable forever.

\begin{itemize}
\tightlist
\item
  \textbf{Gen-2 --- signature-based witness encryption (SWE).} The ``no
  master secret'' property is longevity-optimal for a dead-man's-switch
  (it dissolves W5 and reshapes W1). Adopt when SWE has audited,
  production-grade, cost-acceptable implementations with
  mobile-acceptable ciphertext size.
\item
  \textbf{Layer --- threshold traitor-tracing.} Make a colluding
  committee publicly attributable, closing the \emph{detectability} half
  of W1, when production-ready code exists. This is a layer over the
  base scheme, adoptable independently of gen-1/gen-2.
\item
  \textbf{Gen-3 --- witness encryption from SNARKs.} The trustless
  endgame: the on-chain proof \emph{is} the key, no federation. Tracked
  semi-annually as it matures from theory.
\item
  \textbf{Post-quantum migration.} No deployable PQ threshold-IBE
  exists; harden the \emph{inner} recipient layer first (deployable PQ
  KEMs exist) and track the threshold layer. Keeping
  \texttt{alg}-versioning a hard requirement is the standing mitigation
  for W6.
\item
  \textbf{Tier-2 conditions.} Oracle/API-gated release, opt-in and
  separately tiered, with the data-source trust made explicit (W7).
\item
  \textbf{Full performance evaluation} and the autonomous
  (watch-and-sign) release mode for stronger liveness (§5.3, §8).
\end{itemize}

The architecture spike that chose gen-1 (DKG threshold-IBE) over the
silent-setup and SWE challengers is re-run semi-annually, because the
maturity of SWE and of WE-from-SNARKs is exactly what will flip the
gen-2/gen-3 timing.

\begin{center}\rule{0.5\linewidth}{0.5pt}\end{center}

\section{11. Ethics and Responsible
Disclosure}\label{ethics-and-responsible-disclosure}

Warden is dual-use infrastructure, and we treat the ethics of that
squarely.

\begin{itemize}
\tightlist
\item
  \textbf{What it protects vs.~what it does not.} Warden protects
  \emph{delivery} (a sealed message reaches its recipients once the
  trigger fires) and \emph{content} (only the recipient can read it). It
  does \textbf{not} protect \emph{metadata}: existence, timing, the
  public condition, and the recipient/sender relationship are public
  on-chain by design (§6.3). We disclose this prominently rather than
  letting an ``encrypted, decentralized'' framing imply anonymity it
  does not provide, because a user who \emph{believes} they are
  invisible when they are not is harmed by our overclaim.
\item
  \textbf{No anonymity / press-freedom framing discipline.} The
  press-freedom use case is framed strictly as \emph{work survival}
  (``your work survives you'' --- accident, illness, inability to
  deliver), never as a tool for adversarial-state anonymity, consistent
  with the protocol's honesty guardrail.
\item
  \textbf{Pre-audit status.} The system is testnet and pre-audit; the
  current five-node federation is a single trust domain and so has no
  timing-security, and independent-operator commitments are being sought
  but not yet bound (§9, §10). We frame this work as \emph{design +
  threat model + early deployment} precisely so that publishing it
  before the audit is responsible --- openness supports the audit gate
  rather than substituting for it, and binding founding-operator
  commitments wait on the audit.
\item
  \textbf{Why publishing gives nothing away.} Openness is part of
  Warden's trust model (no serious operator runs a black box), the
  codebase is MIT public-good, and the cryptographic substrate is
  already public (drand). Publishing the design, threat model, and
  deployment experience therefore costs no security; it \emph{is} the
  federation-recruitment vehicle.
\end{itemize}

\begin{center}\rule{0.5\linewidth}{0.5pt}\end{center}

\section{12. Conclusion}\label{conclusion}

Warden takes the deployed timelock idea --- a threshold-BLS federation
that is, unknowingly, an IBE key generator --- and changes exactly one
thing: the identity is the hash of a \emph{resettable on-chain
condition} rather than a future round number, and release is gated on
observing that condition true at finalized depth rather than on a clock.
Around that one change it builds an honest systems story: a double-wrap
envelope that confines the federation to timing and never to content;
resharing that gives the master key permanence across decades of
operator churn; a finality policy for the security-critical chain read
that drand never needed; and a multi-platform implementation running
end-to-end as a five-node testnet federation on Base Sepolia, integrated
with the Maktub protocol (a single trust domain today; independent
operators are being recruited). The cryptography is reused and audited;
the contribution is the system, the resettable-condition threat model,
and a limitation analysis (W1--W11) calibrated to claim exactly the two
properties Warden keeps --- guaranteed-but-deferred delivery and content
confidentiality --- and nothing more. The honest gap between this and
the trustless ideal (witness encryption) is the roadmap, and the
\texttt{alg}-versioned envelope is the bridge across it.

\begin{center}\rule{0.5\linewidth}{0.5pt}\end{center}

\section{\texorpdfstring{Appendix A --- Envelope byte-format
(\texttt{warden-v1})}{Appendix A --- Envelope byte-format (warden-v1)}}\label{appendix-a-envelope-byte-format-warden-v1}

The wire form as implemented (\texttt{core/src/envelope.rs}); blobs are
hex of the compressed-canonical / \texttt{nonce\ ‖\ ct} bytes.

\begin{Shaded}
\begin{Highlighting}[]
\NormalTok{\{}
\NormalTok{  "alg": "warden{-}v1",}
\NormalTok{  "network": "\textless{}warden federation / master{-}key id\textgreater{}",   // bound into both AEAD layers}
\NormalTok{  "condition": \{                                       // PUBLIC; hashed into the IBE identity}
\NormalTok{    "type": "contract",}
\NormalTok{    "chain": 84532,                                    // Base Sepolia (testnet)}
\NormalTok{    "address": "0x\textless{}MaktubCore\textgreater{}",}
\NormalTok{    "fn": "getHeartbeat(uint256)",                     // no executed() getter…}
\NormalTok{    "args": ["\textless{}beatId\textgreater{}"],}
\NormalTok{    "word": 7,                                         // …executed is return field 7 (see §4.2)}
\NormalTok{    "test": \{ "cmp": "==", "value": true \},}
\NormalTok{    "meta": \{ "finality": 32, "tier": 1 \}}
\NormalTok{  \},}
\NormalTok{  "outer": \{                  // condition gate (hybrid IBE — see §4.1)}
\NormalTok{    "ibe": "\textless{}hex\textgreater{}",           // IBE(obk, H(condition)) — the (U,V,W) ciphertext, canonical{-}serialized}
\NormalTok{    "seal": "\textless{}hex\textgreater{}"           // nonce ‖ AEAD\_obk(K\_wrapped), K\_wrapped = ECIES(K, recipientPub)}
\NormalTok{  \},}
\NormalTok{  "inner": \{                  // content layer; recipient{-}only. Warden never touches this.}
\NormalTok{    "ct": "\textless{}hex\textgreater{}"             // nonce ‖ AEAD\_K(pad(payload))}
\NormalTok{  \}}
\NormalTok{\}}
\end{Highlighting}
\end{Shaded}

The identity \texttt{H("warden-cond-v1"\ ‖\ jcs(condition))} is
recomputable by any node from the public \texttt{condition}, so the
envelope need not store it (\texttt{word} is omitted from the canonical
form when 0). No recipient metadata is stored --- the recipient is
implicit in who can ECIES-open \texttt{K\_wrapped}.

\section{\texorpdfstring{Appendix B --- Condition schema
(\texttt{cond-v1})}{Appendix B --- Condition schema (cond-v1)}}\label{appendix-b-condition-schema-cond-v1}

One \texttt{type} per condition, plus a mandatory \texttt{meta}:

\begin{itemize}
\tightlist
\item
  \texttt{contract} ---
  \texttt{\{\ chain,\ address,\ fn,\ args,\ word?,\ test:\ \{cmp,\ value\},\ meta\ \}}
\item
  \texttt{block} ---
  \texttt{\{\ chain,\ field:\ "number"\textbar{}"timestamp",\ cmp,\ value,\ meta\ \}}
\item
  \texttt{event} ---
  \texttt{\{\ chain,\ address,\ sig,\ args,\ meta\ \}}
\item
  \texttt{all\ \textbar{}\ any\ \textbar{}\ not\ \textbar{}\ threshold}
  ---
  \texttt{\{\ of:\ {[}...sub-conditions{]},\ k\ (threshold\ only),\ meta\ \}}
\end{itemize}

Rules:
\texttt{cmp\ ∈\ \{\ ==,\ !=,\ \textgreater{}=,\ \textless{}=,\ \textgreater{},\ \textless{}\ \}}
(string); all \texttt{uint256} args/values are decimal strings (never
JSON numbers); \texttt{word} (contract only, default 0, omitted when 0)
selects the 32-byte ABI return word; serialization is RFC 8785 (JCS) ---
a number-vs-string mismatch yields a different identity and silently
breaks decryption. The three non-negotiable constraints are determinism
(finalized reads only), monotonicity/latching (non-monotone predicates
rejected unless anchored), and finality/reorg safety (§3.1, §4.4). For
\texttt{threshold}/compound conditions spanning multiple chains,
deterministic evaluation across heterogeneous finality semantics is
under-specified in v1 and is deferred to future work.

\section{Appendix C --- Parameter
choices}\label{appendix-c-parameter-choices}

\begin{itemize}
\tightlist
\item
  \textbf{Threshold:} target \texttt{t\ ≈\ ⌈2n/3⌉} (BFT-style), minimum
  \texttt{⌊n/2⌋+1}.
\item
  \textbf{Finality:} federation-wide floor (Base: prefer L1-finalized;
  the current implementation reads the \texttt{finalized} tag);
  \texttt{C.meta.finality} may raise it, never lower it.
\item
  \textbf{Curve / scheme:} BLS12-381, Boneh--Franklin IBE over threshold
  BLS, RFC 9380 hash-to-curve, FO-CCA transform; AEAD =
  ChaCha20-Poly1305; recipient layer = ECIES over secp256k1 with
  HKDF/SEC1 shared-info.
\end{itemize}

\begin{center}\rule{0.5\linewidth}{0.5pt}\end{center}

\section{References}\label{references}

\emph{(ePrint identifiers are given where applicable.)}

\begin{itemize}
\tightlist
\item
  D. Boneh, M. Franklin. \emph{Identity-Based Encryption from the Weil
  Pairing.} CRYPTO 2001.
\item
  D. Boneh, B. Lynn, H. Shacham. \emph{Short Signatures from the Weil
  Pairing.} ASIACRYPT 2001.
\item
  A. Boldyreva. \emph{Threshold Signatures, Multisignatures and Blind
  Signatures Based on the Gap-Diffie-Hellman-Group Signature Scheme.}
  PKC 2003.
\item
  A. Shamir. \emph{How to Share a Secret.} CACM 1979.
\item
  P. Feldman. \emph{A Practical Scheme for Non-interactive Verifiable
  Secret Sharing.} FOCS 1987.
\item
  T. Pedersen. \emph{Non-Interactive and Information-Theoretic Secure
  Verifiable Secret Sharing.} CRYPTO 1991.
\item
  R. Gennaro, S. Jarecki, H. Krawczyk, T. Rabin. \emph{Secure
  Distributed Key Generation for Discrete-Log Based Cryptosystems.}
  EUROCRYPT 1999.
\item
  E. Fujisaki, T. Okamoto. \emph{Secure Integration of Asymmetric and
  Symmetric Encryption Schemes.} CRYPTO 1999.
\item
  N. Gailly, K. Melissaris, Y. Romailler. \emph{tlock: Practical
  Timelock Encryption from Threshold BLS.} IACR ePrint 2023/189.
\item
  drand / League of Entropy. Project documentation and the drand v2.0
  post-mortem (2025).
\item
  R. Rivest, A. Shamir, D. Wagner. \emph{Time-lock Puzzles and
  Timed-release Crypto.} 1996.
\item
  B. Wesolowski. \emph{Efficient Verifiable Delay Functions.} EUROCRYPT
  2019. / K. Pietrzak. \emph{Simple Verifiable Delay Functions.} ITCS
  2019.
\item
  S. Garg, C. Gentry, A. Sahai, B. Waters. \emph{Witness Encryption and
  its Applications.} STOC 2013.
\item
  McFly: \emph{Signature-Based Witness Encryption.} IACR ePrint
  2022/433. / Compact SWE, ePrint 2024/1477. / eWEB (next-block SWE),
  ePrint 2025/1064.
\item
  \emph{Threshold Encryption with Silent Setup.} IACR ePrint 2024/263.
\item
  \emph{Silent Threshold Traitor-Tracing (ST3).} IACR ePrint 2025/342;
  ACM CCS 2025.
\item
  \emph{Witness Encryption from SNARKs / KZG.} IACR ePrint 2025/1364
  (CRYPTO 2025).
\item
  Threshold Network. \emph{TACo --- Threshold Access Control:
  conditional threshold decryption} (project documentation).
\item
  Shutter Network. \emph{Threshold encryption for front-running
  protection} (project documentation).
\item
  E. Kokoris-Kogias, E. C. Alp, L. Gasser, P. Jovanovic, E. Syta, B.
  Ford. \emph{Calypso: Private Data Management for Decentralized
  Ledgers.} VLDB 2021.
\item
  M. Campanelli, B. David, H. Khoshakhlagh, A. Konring, J. B. Nielsen.
  \emph{Encryption to the Future: A Paradigm for Sending Secret Messages
  to Future (Anonymous) Committees.} ASIACRYPT 2022.
\item
  RFC 9380 --- \emph{Hashing to Elliptic Curves.} / RFC 8785 ---
  \emph{JSON Canonicalization Scheme (JCS).} / RFC 8439 ---
  \emph{ChaCha20-Poly1305.} / RFC 5869 --- \emph{HKDF.}
\item
  F. Valsorda et al.~\emph{The age file-encryption format}, v1
  (\texttt{age-encryption.org/v1}; C2SP specification).
\item
  Certicom Research. \emph{SEC 1: Elliptic Curve Cryptography, Version
  2.0} (SECG, 2009) --- ECIES and EC primitives.
\item
  Warden project: source repository ---
  https://github.com/bytesbrains/warden (MIT); overview / operator
  information --- https://warden.bytesbrains.com; specification
  (\texttt{docs/}) ---
  https://github.com/bytesbrains/warden/tree/main/docs.
\end{itemize}

\end{document}
